\documentclass[11pt,a4paper]{article}
\usepackage[utf8]{inputenc}
\usepackage[T1]{fontenc}
\usepackage[french]{babel}
\usepackage{amsmath, amssymb}
\usepackage{geometry}
\geometry{hmargin=2cm, vmargin=2cm}

\begin{document}

\begin{center}
    \textbf{GRENOBLE INP Esisar UGA, VALENCE} \\
    \textbf{MT331 Interro. 1 lundi 4 mai 2026 45'} \\
    \textbf{MT331 2025-2026} \\
    Prénom NOM: ...................................
\end{center}

\vspace{0.5cm}

\section*{Exercice 1}
Répondre au bout de la ligne sans justification. $|A|$ désigne le cardinal de l'ensemble A.

\noindent \textbf{1.} Soit $k$ un entier naturel et $A$ et $B$ deux ensembles avec $|A|=n$ et $|B|=p$.
\begin{enumerate}
    \item[(a)] le cardinal de $A \times B$ est : $n \times p$
    % Raisonnement : Par définition, le cardinal du produit cartésien de deux ensembles finis 
    % est égal au produit des cardinaux de ces deux ensembles.
    
    \item[(b)] le cardinal de $A^k$ est : $n^k$
    % Raisonnement : A^k désigne l'ensemble des applications d'un ensemble à k éléments vers A. 
    % Pour chacun des k éléments de l'ensemble de départ, on a n choix possibles dans A.
    
    \item[(c)] le nombre de k-listes d'éléments de A est : $n^k$
    % Raisonnement : Une k-liste (ou k-uplet) est une suite ordonnée de k éléments où les répétitions 
    % sont autorisées. C'est mathématiquement équivalent à un élément de A^k, d'où n^k possibilités.
    
    \item[(d)] si $k \le n$, le nombre de k-listes sans répétition d'éléments de A est : $\frac{n!}{(n-k)!}$
    % Raisonnement : Il s'agit du nombre d'arrangements de k éléments parmi n (noté A_n^k). 
    % On a n choix pour le premier, (n-1) pour le second... jusqu'à (n-k+1) pour le k-ième.
    
    \item[(e)] le nombre de façons différentes d'ordonner les n éléments de A est : $n!$
    % Raisonnement : Ordonner n éléments distincts revient à chercher le nombre de permutations de cet ensemble. 
    % Le résultat est factorielle n (n!).
    
    \item[(f)] si $k \le n$, le nombre de parties à k éléments de A est : $\binom{n}{k} = \frac{n!}{k!(n-k)!}$
    % Raisonnement : Choisir un sous-ensemble (une partie) de k éléments parmi n revient à calculer 
    % le nombre de combinaisons. Ici, contrairement aux listes, l'ordre n'a pas d'importance.
    
    \item[(g)] le nombre de parties de A est : $2^n$
    % Raisonnement : L'ensemble de toutes les parties de A (noté P(A)) a pour cardinal 2^{|A|}. 
    % On peut aussi comprendre cela en disant que pour construire une partie, chaque élément de A 
    % a exactement 2 choix : appartenir ou ne pas appartenir à cette partie.
\end{enumerate}

\noindent \textbf{2.} Soit $(\Omega, \mathcal{P}(\Omega), \mathbb{P})$ un espace probabilisé fini, et A et B deux événements de $\Omega$.
\begin{enumerate}
    \item[(a)] $\mathbb{P}(\Omega) = 1$
    % Raisonnement : C'est l'axiome fondamental de la définition d'une probabilité : la probabilité 
    % de l'univers entier (l'événement certain) est égale à 1.
    
    \item[(b)] $\mathbb{P}(\emptyset) = 0$
    % Raisonnement : Propriété directe découlant des axiomes : l'événement impossible a une probabilité nulle.
    
    \item[(c)] Si $A \cap B = \emptyset$, alors $\mathbb{P}(A \cup B) = \mathbb{P}(A) + \mathbb{P}(B)$
    % Raisonnement : C'est l'axiome d'additivité pour des événements incompatibles (disjoints). 
    % Comme ils ne peuvent pas se produire en même temps, la probabilité de l'un ou de l'autre est la somme des deux.
    
    \item[(d)] $\mathbb{P}(\overline{A}) = 1 - \mathbb{P}(A)$
    % Raisonnement : Un événement et son contraire forment une partition de l'univers Omega. 
    % Donc P(A) + P(A_barre) = P(Omega) = 1.
    
    \item[(e)] $\mathbb{P}(A \cup B) = \mathbb{P}(A) + \mathbb{P}(B) - \mathbb{P}(A \cap B)$
    % Raisonnement : Formule générale de l'additivité (principe d'inclusion-exclusion). On doit soustraire 
    % P(A inter B) car les éléments communs à A et B ont été comptés deux fois (une fois dans P(A) et une fois dans P(B)).
    
    \item[(f)] Par définition, $\mathbb{P}_A(B) = \mathbb{P}(B|A) = \frac{\mathbb{P}(A \cap B)}{\mathbb{P}(A)}$
    % Raisonnement : C'est la définition formelle de la probabilité conditionnelle de B sachant que 
    % l'événement A est déjà réalisé (sous réserve que P(A) > 0).
    
    \item[(g)] Si $(A_i)_{1\le i\le n}$ est un système complet d'événements de $\Omega$, alors $\sum_{i=1}^{n}\mathbb{P}(B|A_i)\mathbb{P}(A_i) = \mathbb{P}(B)$
    % Raisonnement : Il s'agit de la formule des probabilités totales. Le système complet d'événements 
    % partitionne l'univers de manière disjointe, permettant de décomposer le calcul de P(B) selon les différents scénarios A_i.
\end{enumerate}

\vspace{0.5cm}

\section*{Exercice 2}
On souhaite ranger sur une étagère 4 livres de mathématiques, 6 livres de physique, et 3 de chimie. 
De combien de façons peut-on effectuer ce rangement :

\noindent \textbf{1. si les livres doivent être groupés par matières.} \\
On choisit successivement :
\begin{itemize}
    \item l'ordre des 3 groupes : $3!$ possibilités
    % Raisonnement : On a 3 blocs distincts à agencer sur l'étagère (le bloc M, le bloc P et le bloc C). 
    % Il y a 3! façons d'ordonner ces 3 blocs (ex: M-P-C, C-M-P, etc.).
    
    \item l'ordre des livres de maths : $4!$ possibilités
    % Raisonnement : À l'intérieur du bloc de mathématiques, on doit ranger les 4 livres distincts (4! possibilités).
    
    \item l'ordre des livres de physique : $6!$ possibilités
    % Raisonnement : À l'intérieur du bloc de physique, on fait de même avec les 6 livres (6! possibilités).
    
    \item l'ordre des livres de chimie : $3!$ possibilités
    % Raisonnement : À l'intérieur du bloc de chimie, on ordonne les 3 livres (3! possibilités).
\end{itemize}
Donc (principe multiplicatif) : $3! \times 4! \times 6! \times 3!$
% Raisonnement global : Comme ces choix s'effectuent de manière successive et indépendante, 
% le principe multiplicatif stipule qu'on multiplie le nombre de possibilités de chaque étape.

\vspace{0.3cm}
\noindent \textbf{2. si seuls les livres de mathématiques doivent être groupés.} \\
On choisit successivement :
\begin{itemize}
    \item le nombre de livres qui précèdent le bloc de maths : 10 possibilités
    % Raisonnement : Il y a au total 6 + 3 = 9 livres qui ne sont pas des maths. Le bloc indissociable des 4 livres 
    % de maths peut se positionner à 10 endroits différents : soit avant le premier livre non-maths, 
    % soit entre le 1er et le 2e... soit tout à la fin après le 9e livre non-maths.
    
    \item l'ordre des livres de maths : $4!$
    % Raisonnement : On ordonne les 4 livres de mathématiques au sein de leur propre bloc.
    
    \item l'ordre des livres non maths : $9!$
    % Raisonnement : On ordonne les 9 autres livres (physique et chimie mélangés) sur le reste de l'étagère.
\end{itemize}
Donc avec le principe multiplicatif : $10 \times 4! \times 9! = 10! \times 4!$
% Autre raisonnement très élégant : On peut considérer le bloc de maths comme un seul gros "super-livre". 
% On a alors à ranger 9 livres isolés + 1 super-livre = 10 éléments. Il y a 10! façons de les ordonner. 
% Ensuite, on multiplie par les 4! façons de permuter les livres de maths à l'intérieur de leur super-livre.

\vspace{0.5cm}

\section*{Exercice 3}
On considère les lettres du mot EVENEMENT (sans accents). Combien d'anagrammes différents peut-on former ? \\
% Raisonnement préliminaire : Le mot contient 9 lettres en tout. Cependant, certaines lettres se répètent : 
% la lettre E apparaît 4 fois, et la lettre N apparaît 2 fois. Les lettres V, M, T sont uniques.
On choisit successivement :
\begin{itemize}
    \item la position des E : $\binom{9}{4}$
    % Raisonnement : Parmi les 9 cases vides de notre anagramme, on choisit les 4 emplacements où seront placés les 'E'.
    
    \item la position des N : $\binom{5}{2}$
    % Raisonnement : Il reste 9 - 4 = 5 cases vides. On choisit 2 emplacements pour y placer les 'N'.
    
    \item l'ordre pour V, M et T : $3!$
    % Raisonnement : Il reste exactement 3 cases pour 3 lettres toutes distinctes (V, M, T). On les permute (3!).
\end{itemize}
soit en tout : $\binom{9}{4} \times \binom{5}{2} \times 3! = \frac{9 \times 8 \times 7 \times 6}{4!} \times \frac{5 \times 4}{2!} \times 3! = \frac{9!}{4! \times 2!}$
% Raisonnement alternatif (formule des permutations avec répétitions) : On calcule le nombre total de permutations (9!) 
% que l'on divise par le produit des factorielles des répétitions pour annuler les doublons visuels : 9! / (4! * 2!).

\vspace{0.3cm}
Parmi ces anagrammes, combien commencent et se terminent par la même lettre ? \\
% Raisonnement : Pour qu'un anagramme commence et finisse par la même lettre, il faut impérativement que cette lettre 
% soit présente au moins deux fois dans le mot de base. Les seuls choix possibles sont donc la lettre 'N' ou la lettre 'E'.
Il y a 2 cas pour la lettre identique :
\begin{enumerate}
    \item Les N : Il reste 7 lettres à placer (dont 4 E), soit $\frac{7!}{4!}$ possibilités
    % Raisonnement : On place un 'N' en première position et un 'N' en dernière position (figés). 
    % Il reste 7 places au milieu à combler avec le reste des lettres : 4 'E', 1 'V', 1 'M', 1 'T'. 
    % Le nombre de structures possibles au milieu est calculé par 7! / 4! (car les 4 'E' sont interchangeables).
    
    \item Les E : Il reste 7 lettres à placer (dont 2 E et 2 N), soit $\frac{7!}{2!2!}$ possibilités
    % Raisonnement : On place un 'E' au début et un 'E' à la fin. Il reste 7 places au centre. 
    % Les lettres restantes à placer sont : 2 'E' restants, 2 'N', 1 'V', 1 'M', 1 'T'. 
    % Le nombre d'anagrammes de ce bloc central est donné par 7! / (2! * 2!).
\end{enumerate}
Donc avec le principe additif (les deux cas étant mutuellement exclusifs) : $\frac{7!}{2^2} + \frac{7!}{4!} = \frac{3 \times 7! + 7!}{12} = \frac{4 \times 7!}{12} = \frac{7!}{3}$

\vspace{0.5cm}

\section*{Exercice 4}
On prend 5 cartes simultanément dans un jeu de 32 cartes. Quelle est la probabilité d'obtenir au moins un as ?

On commencera par modéliser cette expérience aléatoire en donnant l'ensemble $\Omega$ des issues et en donnant la probabilité $\mathbb{P}$ sur $\Omega$. \\
Soit $E$ l'ensemble des 32 cartes, $\Omega$ est l'ensemble des parties à 5 éléments de $E$, et $\mathbb{P}$ est la probabilité équirépartie.
% Raisonnement de modélisation : Le tirage étant simultané, l'ordre n'a aucune importance et il n'y a pas de remise. 
% L'univers Omega est donc l'ensemble de tous les sous-ensembles (combinaisons) possibles de 5 cartes choisies parmi 32. 
% Le nombre total d'issues est card(Omega) = \binom{32}{5}. Le tirage s'effectuant de manière juste, on applique la loi d'équiprobabilité.

En notant :
\begin{itemize}
    \item $A$ : obtenir au moins un as
    \item $\overline{A}$ : ne pas obtenir d'as
\end{itemize}
on a :
$$ \mathbb{P}(\overline{A}) = \frac{|\overline{A}|}{|\Omega|} = \frac{\binom{28}{5}}{\binom{32}{5}} $$
% Raisonnement : Calculer la probabilité d'avoir "au moins un as" directement demanderait d'additionner les cas : 
% avoir exactement 1 as, 2 as, 3 as, ou 4 as. Il est beaucoup plus rapide de passer par l'événement contraire "n'obtenir aucun as".
% Un jeu de 32 cartes contient 4 as. Les cartes qui ne sont pas des as sont au nombre de 32 - 4 = 28. 
% Pour n'avoir aucun as, il faut choisir nos 5 cartes uniquement parmi ces 28 cartes non-as, d'où \binom{28}{5} possibilités.

Donc :
$$ \mathbb{P}(A) = 1 - \mathbb{P}(\overline{A}) = 1 - \frac{28 \times 27 \times 26 \times 25 \times 24}{32 \times 31 \times 30 \times 29 \times 28} $$
% Raisonnement final : On applique la formule fondamentale des probabilités complémentaires : P(A) = 1 - P(A_barre).

\end{document}